\section*{一、选题背景与意义}
\addcontentsline{toc}{section}{一、选题背景与意义}
\subsection*{1.1 研究背景}
\addcontentsline{toc}{subsection}{1.1 研究背景}
TCP/IP 协议栈是互联网的基石,但其设计之初对安全性考虑不足,且经过数十年的演进与打补丁,协议实现内部及协议之间形成了大量复杂的跨层交互。近十年来的研究表明,这些跨层交互中的歧义(ambiguity)持续地催生新的\textbf{链外(off-path)攻击}:攻击者无需位于通信路径上、无需窃听报文,仅凭侧信道即可推断 TCP 连接的关键状态(如序列号),进而劫持、重置或污染连接\supercite{feng2025cacm}。

这一攻击面呈现出\textbf{持续再生}的特征。自 2016 年 Cao 等人揭示 RFC 5961 引入的全局 challenge-ACK 速率限制可被用作侧信道(CVE-2016-5696)以来\supercite{cao2016offpath},学术界沿"IP–ICMP–TCP 跨层交互"这一主线不断发现新漏洞:混合/降级的 IPID 分配侧信道\supercite{feng2020mixedipid}、绕过 PMTUD 的 IP 分片攻击\supercite{feng2022pmtud}、NAT 环境下的 TCP 劫持\supercite{yang2024natwifi}、Wi-Fi 加密帧的帧长侧信道\supercite{wang2025packetsize},乃至作者所在课题组在 2025 年 CCS 上发表的、利用路径 MTU 发现(PMTUD)打破 IP 地址共享场景下连接隔离的最新攻击\supercite{feng2025pmtudisolation}。与此同时,DNS 等无连接协议同样不断出现侧信道投毒攻击\supercite{man2020saddns},基于 IP 伪造的 TCP 载荷可靠注入也被证明可行\supercite{pan2024tcpspoofing}。

一个尤为值得警惕的规律是:\textbf{许多新攻击面恰恰由防御机制本身引入}——challenge-ACK 限速是为抵御盲注入而设,却成了侧信道\supercite{cao2016offpath};源端口随机化\supercite{rfc6056}是为抵御 off-path 注入而设,却被侧信道反推\supercite{man2020saddns};PMTUD 是为优化传输而设,却被用于打破连接隔离\supercite{feng2025pmtudisolation}。这说明协议安全不是"修补一次即告终结"的静态问题,而是随部署环境(公网 Wi-Fi、VPN、5G CGNAT 等 IP 共享场景)与防御补丁不断演化的\textbf{动态博弈}。

面对层出不穷的协议攻击,\textbf{现有防御呈现三个结构性缺陷}:

\begin{enumerate}[leftmargin=2em,itemsep=1pt,topsep=2pt]
  \item \textbf{零散}。防御以"一种攻击对应一个补丁"的方式分散在标准(RFC 5961\supercite{rfc5961}、RFC 6056\supercite{rfc6056}、BCP38\supercite{rfc2827})、内核实现与中间盒中,缺乏统一的、可持续演进的防御载体;
  \item \textbf{滞后}。检测侧虽已出现 SCENT\supercite{cao2019scent}、SCAD\supercite{man2024scad} 等系统化挖掘协议侧信道的工具,但其本质仍是对已成形漏洞类别的\textbf{被动追赶};从漏洞披露到内核缓解上线往往周期漫长;
  \item \textbf{依赖人工编写内核代码}。具体缓解几乎都需安全专家手工设计、提交并等待上游合入。人工编写内核防御面临\textbf{响应慢、易出错(跨层逻辑复杂、补丁本身可能成为新攻击面)、知识不可积累}三重困境。
\end{enumerate}
与此同时,内核可编程技术 \textbf{eBPF} 的成熟为运行时防御提供了理想载体。eBPF 允许在不修改内核源码、不重启系统的前提下,将自定义程序安全地注入内核热点路径(XDP/tc/kprobe/LSM 等),已被工业界广泛用于高性能负载均衡\supercite{shirokov2018katran}、DDoS 缓解\supercite{fabre2018l4drop}、可扩展防火墙\supercite{miano2019bpfiptables}与云原生运行时安全\supercite{graf2020cilium,tetragon2022,falco2018}。其内置的\textbf{验证器(verifier)} 能在加载前对程序进行形式化的内存安全与终止性证明\supercite{gershuni2019prevail,sun2024stateembedding},构成了"让不可信代码安全运行于内核"的可信基础。

另一条技术线索来自人工智能:\textbf{大模型编码智能体(LLM Coding Agent)} 已展现出自主编写、调试与迭代代码的能力,并在程序搜索与算法发现上取得突破——FunSearch 在程序空间做进化搜索发现了更优算法\supercite{romeraparedes2024funsearch},AlphaEvolve 进一步在系统级代码(数据中心调度、矩阵乘法)上做出可验证改进\supercite{novikov2025alphaevolve};在安全领域,编码智能体已能自主发现真实软件漏洞\supercite{glazunov2024naptime}、完成渗透测试\supercite{deng2024pentestgpt}并利用真实漏洞\supercite{fang2024oneday}。尤为关键的是,已有工作 \textbf{Kgent/KEN} 证明大模型能够从自然语言合成 eBPF 程序并结合符号执行保证语义正确\supercite{yang2024kgent},这为"智能体编写内核防御"提供了直接的可行性证据。

\subsection*{1.2 研究意义}
\addcontentsline{toc}{subsection}{1.2 研究意义}
综合上述背景,本课题提出并构建\textbf{一个由大模型编码智能体自主维护、跑在内核的自进化 eBPF 网络入侵防御系统}。其意义体现在三个层面:

\begin{itemize}[leftmargin=2em,itemsep=1pt,topsep=2pt]
  \item \textbf{理论意义}:针对"防御能力如何持续学习"这一根本问题,本课题提出\textbf{防御演化(Defensive Evolution)} 范式——将防御能力沉淀为"经验证的内核代码"而非神经网络权重,以编码智能体的非梯度迭代替代反向传播。这一范式受"超越梯度的学习(Learning Beyond Gradients)"\supercite{weng2026hl} 启发,有望规避深度学习在持续学习中的灾难性遗忘问题,并为"代码即可持续进化的知识载体"提供一个安全关键(security-critical)领域的实证。
  \item \textbf{现实/工程意义}:把作者及课题组在 off-path 协议攻击上积累的\textbf{攻击者视角}\supercite{feng2025cacm,feng2025pmtudisolation} 反向用于防御,使防御系统能够"主动想象攻击面"并提前生成缓解;以 eBPF 作为可热加载、可验证、贴近协议栈的内核防御载体,有望将协议防御从"被动打补丁"转变为"主动、自动、持续的免疫",显著缩短从新攻击出现到防御上线的周期。
  \item \textbf{学科交叉意义}:本课题处在\textbf{网络协议安全、操作系统内核(eBPF)、大模型智能体}三个方向的交叉点上,既是对"AI for Security"的具体落地,也为"AI 自主维护安全关键系统"的可靠性(verifier + 沙箱 + 自动回滚)提供方法论探索。
\end{itemize}
\section*{二、国内外研究现状综述}
\addcontentsline{toc}{section}{二、国内外研究现状综述}
本节从四条线索梳理国内外研究现状,并在每条线索末尾指出其局限,从而引出本课题的切入点。

\subsection*{2.1 链外协议攻击及其防御:攻防同源、防御滞后}
\addcontentsline{toc}{subsection}{2.1 链外协议攻击及其防御:攻防同源、防御滞后}
\textbf{攻击侧},off-path TCP/IP 攻击自 Gilad 等确立"仅能伪造源地址、不能窃听"的威胁模型以来\supercite{gilad2012offpath},已发展为协议安全的活跃方向。清华大学徐恪、李琦、冯学伟等所在的网络体系结构与安全课题组,沿"跨层交互歧义"主线连续产出成果:混合 IPID 分配侧信道\supercite{feng2020mixedipid}、降级 IPID 劫持、PMTUD 分片攻击\supercite{feng2022pmtud}、ICMP 重定向攻击的"复活"、WPA 环境下的中间人攻击、NAT 网络 TCP 劫持\supercite{yang2024natwifi}、Wi-Fi 帧长侧信道\supercite{wang2025packetsize},并以 CACM 综述将其统一为"TCP/IP 跨层漏洞"框架\supercite{feng2025cacm}。作者本人参与的 CCS'25 工作进一步揭示:即便是被普遍信任的 PMTUD 机制,在公网 Wi-Fi/VPN/5G 等 IP 地址共享场景下仍会打破连接隔离,使攻击者平均 220 秒、70\% 成功率地推断受害连接序列号,50 个真实网络中 38 个脆弱\supercite{feng2025pmtudisolation}。UCR 的 Qian 团队则在 DNS 方向揭示了 SAD DNS 等侧信道投毒\supercite{man2020saddns},CISPA 的 Pan 与 Rossow 证明了伪造握手后仍可可靠注入 TCP 载荷\supercite{pan2024tcpspoofing}。

\textbf{防御侧},标准化手段(RFC 5961 加固\supercite{rfc5961}、源端口随机化\supercite{rfc6056}、入口过滤 BCP38\supercite{rfc2827})多为"建议非强制",部署参差,且如前所述常被转化为新的攻击面。系统化检测方面,Cao 等提出 SCENT 将 TCP 侧信道归结为连接间"非干扰性"的违反并半自动挖掘漏洞\supercite{cao2019scent},Man 等的 SCAD 将其推广为通用自动检测\supercite{man2024scad};课题组自身也在无连接协议信息泄露上提出了 StateShield 实时防御\supercite{li2025stateshield}。

\textbf{局限与切入点}:现有防御本质是\textbf{逐漏洞、逐协议、人工驱动}的;检测工具能"发现"漏洞却不能"自动修复并部署";防御知识散落在论文与补丁里无法自动积累复用。这正是本课题要以"自主进化的内核防御系统"填补的缺口。

\subsection*{2.2 eBPF 网络安全与运行时防御:载体成熟、智能化不足}
\addcontentsline{toc}{subsection}{2.2 eBPF 网络安全与运行时防御:载体成熟、智能化不足}
eBPF 已成为内核可编程与运行时安全的事实标准载体。工业界以 Katran\supercite{shirokov2018katran}、Cloudflare L4Drop\supercite{fabre2018l4drop} 在 XDP 层实现高性能负载均衡与 DDoS 缓解;学术界以 bpf-iptables\supercite{miano2019bpfiptables} 复刻并加速防火墙语义;Cilium\supercite{graf2020cilium}、Falco\supercite{falco2018}、Tetragon\supercite{tetragon2022} 则把 eBPF 用于云原生网络策略与运行时安全可观测/强制执行(Tetragon 可在系统调用完成前同步阻断)。在安全保障方面,eBPF 验证器是关键:PREVAIL 以抽象解释支持任意循环验证\supercite{gershuni2019prevail},K2 以程序综合生成"安全且高效"的内核扩展\supercite{xu2021k2},Sun 等通过状态嵌入检测验证器自身缺陷\supercite{sun2024stateembedding}。同时,eBPF 也被研究为攻击面——跨容器攻击\supercite{he2023crosscontainer}、Evil Packet Filter\supercite{jin2023epf} 揭示了其武器化风险,反向论证了"eBPF 防御库本身需要被安全地维护"。

\textbf{局限与切入点}:上述系统的防御策略\textbf{几乎全部由人工编写};规则/程序的演进依赖工程师手工迭代。当攻击快速演化时,人工维护成为瓶颈。eBPF 提供了优秀的"载体"与"安全闸门",但缺少一个\textbf{自主的"维护者"}——这正是引入编码智能体的动机。

\subsection*{2.3 大模型编码智能体用于安全:能力涌现、面向防御者尚少}
\addcontentsline{toc}{subsection}{2.3 大模型编码智能体用于安全:能力涌现、面向防御者尚少}
大模型在安全领域的应用近两年迅速成熟。进攻侧,PentestGPT 实现自主渗透\supercite{deng2024pentestgpt},LLM 智能体可自主利用真实一日漏洞\supercite{fang2024oneday},Google 的 Naptime/Big Sleep 让智能体在真实软件中发现此前未知的内存安全漏洞\supercite{glazunov2024naptime};评测侧出现了 Cybench\supercite{zhang2025cybench}、NYU CTF Bench\supercite{shao2024nyuctf} 等基准。代码维护侧,自动程序修复(APR)综述显示 LLM 已成为该领域主流方法\supercite{zhang2024aprsurvey}。与本课题最直接相关的是 \textbf{Kgent/KEN}:它用大模型从自然语言合成 eBPF 程序,并以符号执行 + 反馈回路保证语义等价,正确率达 80\%\supercite{yang2024kgent}。

在工程载体层面,\textbf{通用编码智能体}已走向成熟并产品化:Anthropic 的 Claude Code\supercite{anthropic2025claudecode}、Anysphere 的 Cursor\supercite{cursor2024}、OpenAI 的 Codex CLI\supercite{openai2025codexcli}、Cognition 的 Devin\supercite{cognition2024devin},以及开源的 SWE-agent\supercite{yang2024sweagent}、OpenHands\supercite{wang2024openhands} 等,均能在仓库级任务上自主完成"编辑—运行—调试"的闭环,并依托标准化的智能体—计算机接口(agent–computer interface)调用工具。然而它们面向\textbf{通用软件工程},并不内置内核 eBPF 防御所需的领域能力——内核 verifier 报错的结构化解析、内核态隔离沙箱、攻防靶场、防御档案与探索信号;直接套用通用智能体难以胜任本任务。这构成本课题的一项关键设计:与其直接调用通用智能体,不如把上述领域工具与反馈固化为面向防御演化的\textbf{专用编码智能体 harness}(domain-specific agent scaffold)。

\textbf{局限与切入点}:(1) 现有"LLM 写 eBPF"工作(如 Kgent\supercite{yang2024kgent})是\textbf{一次性、按需合成},而非\textbf{长期自主维护一整套防御库 + 运行时反馈 + 持续进化};(2) 安全应用多集中在\textbf{进攻侧},面向\textbf{防御者}、面向\textbf{协议攻击}的自主防御研究尚少;(3) 研究普遍指出 LLM 生成代码的可靠性存疑(不安全代码比例可达 9.8\%–42.1\%\supercite{seccodegen2025}),因此"让 AI 写内核代码"必须配套强制的验证与沙箱——这恰是 eBPF verifier 的独特价值所在;(4) 通用编码智能体\supercite{yang2024sweagent,anthropic2025claudecode}缺少内核/eBPF 的领域工具与反馈闭环,需要为本任务设计\textbf{专用 harness}。

\subsection*{2.4 启发式/进化式学习与自进化智能体:方法论根基}
\addcontentsline{toc}{subsection}{2.4 启发式/进化式学习与自进化智能体:方法论根基}
本课题"非梯度、持续进化"的方法论有坚实的文献支撑。其直接灵感是"超越梯度的学习"\supercite{weng2026hl}:主张编码智能体可通过\textbf{持续编辑一个软件系统}(而非更新权重)实现"学习",具备可解释、样本高效、可回归测试、抗灾难性遗忘等优势。在程序搜索层面,FunSearch\supercite{romeraparedes2024funsearch}、AlphaEvolve\supercite{novikov2025alphaevolve}、EoH\supercite{liu2024eoh} 证明"冻结的大模型作进化算子 + 自动评估器作选择压力"能在程序/系统级代码上做出可验证改进。在自我改进层面,Reflexion\supercite{shinn2023reflexion} 以语言化反思替代梯度积累经验,Self-Debug\supercite{chen2024selfdebug} 让模型据执行/编译错误自我修正,Darwin–Gödel Machine\supercite{zhang2025dgm} 与 Voyager\supercite{wang2023voyager} 展示了带档案库、可持续自我进化的智能体。在探索机制层面,基于预测误差与新颖性的探索方法 ICM\supercite{pathak2017icm}、随机网络蒸馏 RND\supercite{burda2019rnd} 提供了"新颖性/预测误差"驱动的探索信号,新颖性搜索\supercite{lehman2011novelty} 与 MAP-Elites\supercite{mouret2015mapelites} 提供了"维护多样化优质解集"的算法范式,其思想可追溯至 Schmidhuber 关于内在激励与新颖性的早期理论\supercite{schmidhuber2010creativity}。

在自主网络防御(ACD)方向,自主计算\supercite{kephart2003autonomic} 与自保护软件系统\supercite{yuan2014selfprotecting} 奠定了 MAPE-K 自适应闭环的理论基础,CybORG/CAGE\supercite{standen2021cyborg} 提供了 RL 防御基准,近期亦有将大模型用作自主防御者的探索\supercite{castro2025llmacd}。

\textbf{局限与切入点}:这些方法论尚未被系统性地应用于"内核网络防御"这一安全关键、强可验证、反馈廉价的场景。本课题正是要把"新颖性与覆盖度引导的主动探索 + LLM 进化式程序合成 + verifier/沙箱选择压力 + 新颖性档案"整合为面向 eBPF 防御的\textbf{防御演化闭环}。

\section*{三、研究内容与拟达到的目标}
\addcontentsline{toc}{section}{三、研究内容与拟达到的目标}
\subsection*{3.1 总体目标}
\addcontentsline{toc}{subsection}{3.1 总体目标}
设计并实现一个\textbf{自进化 eBPF 网络入侵防御系统}:由大模型编码智能体在"防御演化"范式下,自主地发现网络协议攻击面、合成并迭代 eBPF 防御程序、经验证与沙箱评估后热加载到内核数据面,形成一个\textbf{持续生长、可审计、形式化可验证}的防御代码库,在保持低误报与线速性能的同时,逐步提升对已知与未知协议攻击的防御能力。

\subsection*{3.2 主要研究内容}
\addcontentsline{toc}{subsection}{3.2 主要研究内容}
\textbf{研究内容一:面向 eBPF 防御的完整专用 harness。} 将"由编码智能体生成 eBPF 防御并保证其可用、安全"的全过程,封装为一套完整的领域专用 harness:涵盖防御程序的生成、编译、内核 verifier 校验、隔离沙箱测试、加载与运行时反馈接口;并解决两个关键工程问题——生成质量(智能体生成的 eBPF 能编译通过、通过 verifier,且确实拦截攻击而不误伤正常流量)与运行安全(以编译/verifier 报错驱动自我修正,异常时自动回滚,保证不致网络中断或误伤正常流量)。这区别于直接调用 Claude Code、Cursor 等通用编码智能体\supercite{yang2024sweagent,anthropic2025claudecode,cursor2024}。

\textbf{研究内容二:面向真实协议攻击的防御生成与评测。} 复用 PMTUD/NAT/IPID 等真实协议攻击\supercite{feng2025pmtudisolation,yang2024natwifi,feng2020mixedipid}搭建攻防靶场,让系统针对这些攻击自动生成并迭代防御;量化检测率、误报率、XDP 吞吐与时延开销,以及达到目标防御所需的人工介入量,并与人工编写规则、Kgent\supercite{yang2024kgent}、Tetragon\supercite{tetragon2022}/Falco\supercite{falco2018} 等对比;对未让系统接触过的新攻击单独测试,考察其泛化能力。

\subsection*{3.3 拟达到的目标}
\addcontentsline{toc}{subsection}{3.3 拟达到的目标}
\begin{itemize}[leftmargin=2em,itemsep=1pt,topsep=2pt]
  \item 完成系统原型与智能体工具链闭环(在 PacketScope/Guarder 上落地);
  \item 在典型协议攻击集合(含课题组自有攻击)上,系统在\textbf{无人工编写规则}前提下自主生成有效 eBPF 防御,关键攻击的防御成功率与人工专家规则相当或更优,误报率控制在可用范围;
  \item 验证"自进化"特性:面对\textbf{预留的未知攻击},系统能在无人工补丁介入下经若干轮迭代自主提升防御覆盖;
  \item XDP 数据面在开启防御时的吞吐损失与时延开销维持在工程可接受水平;
  \item 形成一套可运行、可复现的开源系统与攻防评测靶场。
\end{itemize}
\section*{四、拟解决的关键问题}
\addcontentsline{toc}{section}{四、拟解决的关键问题}
\begin{enumerate}[leftmargin=2em,itemsep=1pt,topsep=2pt]
  \item \textbf{可靠性——如何让智能体自动生成的内核代码安全可控?} 内核态 eBPF 一旦出错可致网络中断或误伤全部流量,而 LLM 生成代码可靠性存疑\supercite{seccodegen2025}。关键在于把\textbf{现成的内核安全设施}(内核 verifier、静态沙箱、隔离环境对抗验证、自动回滚)\textbf{接入智能体闭环}:以编译/verifier 报错驱动自修正,沙箱重放把关,异常即回滚——本文不重造这些基础设施,而是组织它们构成多层可靠闭环。
  \item \textbf{生成质量——智能体生成的 eBPF 怎样才"能用"?} 它必须能编译通过、通过内核 verifier、且确实拦截目标攻击而不误伤正常流量。关键在于把编译/verifier 报错与沙箱重放结果作为结构化反馈,驱动智能体迭代修正,直到产出可用且达标的防御。
  \item \textbf{评测——如何证明系统"真的有用"?} 不仅要防住已见过的攻击,对未让系统接触过的新攻击也要能防。关键在于构建真实攻击靶场,给出检出率、误报率、性能开销与人工介入量等可复现、可对比的度量,并与人工规则及现有工具对比。
\end{enumerate}
\section*{五、研究方法与技术路线}
\addcontentsline{toc}{section}{五、研究方法与技术路线}
% 系统体系结构图(TikZ)。依赖主文件已定义颜色 thupurple/thublue 与 tikz 库。
\begin{figure}[htbp]
  \centering
  \resizebox{\textwidth}{!}{%
  \begin{tikzpicture}[font=\small, node distance=0.62cm and 0.7cm,
     box/.style={rectangle, rounded corners=3pt, draw=thupurple, line width=0.8pt,
                 fill=thupurple!8, align=center, text width=2.55cm, minimum height=1.05cm, inner sep=3pt},
     wide/.style={rectangle, rounded corners=3pt, draw=thupurple, line width=0.8pt,
                 fill=thupurple!8, align=center, text width=6.4cm, minimum height=1.05cm},
     side/.style={rectangle, rounded corners=3pt, draw=thublue, line width=0.8pt,
                 fill=thublue!8, align=center, text width=3.0cm, minimum height=1.1cm},
     lbl/.style={text=thupurple!85, font=\footnotesize\bfseries, align=center, rotate=90},
     arr/.style={-{Stealth[length=2.6mm]}, line width=0.9pt, draw=thupurple!85},
     barr/.style={-{Stealth[length=2.6mm]}, line width=0.9pt, draw=thublue!80, dashed}]
   % ===== 智能体层 =====
   \node[box] (synth) {进化式 eBPF 合成器\\[1pt]\scriptsize LLM 生成/改写 C 代码};
   \node[box, left=of synth] (explore) {攻击面探索器\\[1pt]\scriptsize 新颖性/覆盖度};
   \node[box, right=of synth] (reflect) {反思迭代器\\[1pt]\scriptsize Reflexion/Self-Debug};
   \node[lbl, left=0.35cm of explore] {智能体层};
   % ===== 可信加载闭环层 =====
   \node[box, below=1.25cm of explore] (compile) {clang 编译\\[1pt]\scriptsize $\to$ BPF 字节码};
   \node[box, right=of compile] (verifier) {内核 verifier\\[1pt]\scriptsize 内存安全·有界·终止};
   \node[box, right=of verifier] (sandbox) {隔离沙箱\\[1pt]\scriptsize 攻击流量对抗评估};
   \node[box, right=of sandbox] (rollback) {自动回滚\\[1pt]\scriptsize 异常即撤回};
   \node[lbl, left=0.35cm of compile] {可信闭环};
   % ===== 内核数据面层 =====
   \node[wide, below=1.25cm of compile.south west, anchor=north west] (tail)
        {内核数据面(XDP):tail-call 防御链\\[1pt]\scriptsize 过滤器 $f_0 \to f_1 \to \cdots \to$ \texttt{XDP\_PASS}\;/\;\texttt{XDP\_DROP}};
   \node[box, right=of tail] (tele) {连接追踪\\与遥测\\[1pt]\scriptsize 检出/误报/时延};
   \node[lbl, left=0.35cm of tail] {内核 XDP};
   % ===== 右侧:防御库 + 靶场 =====
   \node[side, right=1.0cm of reflect] (archive) {防御库\\(MAP-Elites 档案)\\[1pt]\scriptsize 多样化精英防御};
   \node[side, below=2.4cm of archive] (range) {攻防靶场 / 红队\\[1pt]\scriptsize 课题组 PMTUD/NAT/IPID 攻击};
   % ===== 数据流 =====
   \draw[arr] (explore) -- (synth);
   \draw[arr] (synth) -- (reflect);
   \draw[arr] (synth.south) -- (compile.north) node[midway,right=1pt,font=\scriptsize,text=thupurple!75]{候选 eBPF};
   \draw[arr] (compile) -- (verifier);
   \draw[arr] (verifier) -- (sandbox);
   \draw[arr] (sandbox) -- (rollback);
   \draw[arr] (sandbox.south) -- (tail.north) node[midway,right=1pt,font=\scriptsize,text=thupurple!75]{热加载};
   % 反馈回路(数据面 -> 智能体,走最左)
   \draw[arr] (tele.east) -- ++(0.5,0) |- (reflect.east);
   \draw[arr, draw=thupurple!70] (tail.west) -- ++(-1.5,0) |- (explore.west)
        node[pos=0.25,left=1pt,font=\scriptsize,text=thupurple!75]{遥测反馈};
   % 档案与合成器互通
   \draw[barr] (archive.west) -- (reflect.east);
   \draw[barr] (synth.north) .. controls +(0,0.8) and +(-1.2,0) .. (archive.north west)
        node[pos=0.6,above,font=\scriptsize,text=thublue!75]{采样父代 / 沉淀精英};
   % 靶场喂攻击样本
   \draw[barr] (range.west) -- ++(-0.5,0) |- (explore.east) node[pos=0.75,above,font=\scriptsize,text=thublue!75]{攻击样本};
  \end{tikzpicture}}
  \caption{自进化 eBPF 网络入侵防御系统总体架构:编码智能体在“探索—合成—可信验证—部署反馈—档案”闭环中维护内核防御库。}
  \label{fig:arch}
\end{figure}

\subsection*{5.1 总体技术路线}
\addcontentsline{toc}{subsection}{5.1 总体技术路线}
本课题采用"\textbf{理论建模 $\to$ 系统实现 $\to$ 攻防评测 $\to$ 迭代改进}"的研究方法,以课题组开源项目 PacketScope 的 Guarder 模块为工程基线(其已具备 XDP 数据面、cilium/ebpf 一键编译、规则热加载、MCP 接入与多维遥测,并已有"Agent 可编程 eBPF + 四层安全模型"的原型),逐步构建完整的防御演化闭环。

\subsection*{5.2 防御演化闭环(核心方法)}
\addcontentsline{toc}{subsection}{5.2 防御演化闭环(核心方法)}
系统以一个开放式进化循环运行,各环节均有方法论依据:

% Defensive-Evolution 闭环图(TikZ)。依赖主文件已定义颜色 thupurple 与 tikz 库。
\begin{figure}[htbp]
  \centering
  \resizebox{0.99\textwidth}{!}{%
  \begin{tikzpicture}[font=\small,
     stage/.style={rectangle, rounded corners=4pt, draw=thupurple, line width=0.9pt,
                   fill=thupurple!8, align=center, text width=2.75cm, minimum height=1.55cm},
     arr/.style={-{Stealth[length=3mm]}, line width=1.0pt, draw=thupurple!85}]
   \node[stage] (s1) {\textbf{①攻击面探索}\\[2pt]\scriptsize 新颖性/覆盖度驱动};
   \node[stage, right=1.05cm of s1] (s2) {\textbf{②防御合成}\\[2pt]\scriptsize LLM 作进化算子};
   \node[stage, right=1.05cm of s2] (s3) {\textbf{③可信验证}\\[2pt]\scriptsize 编译·verifier·沙箱};
   \node[stage, right=1.05cm of s3] (s4) {\textbf{④部署与反馈}\\[2pt]\scriptsize XDP 热加载·遥测};
   \node[stage, right=1.05cm of s4] (s5) {\textbf{⑤档案更新}\\[2pt]\scriptsize 质量\,$\times$\,新颖性};
   \draw[arr] (s1) -- (s2);
   \draw[arr] (s2) -- (s3);
   \draw[arr] (s3) -- (s4);
   \draw[arr] (s4) -- (s5);
   % 内含自修正小回路(验证失败 -> 合成)
   \draw[arr, dashed, draw=thupurple!60] (s3.north) .. controls +(0,0.9) and +(0,0.9) .. (s2.north)
        node[midway, above, font=\scriptsize, text=thupurple!80] {verifier 失败 $\to$ Self-Debug 自修正};
   % 总反馈回路(s5 -> s1)
   \draw[arr] (s5.south) -- ++(0,-1.0)
        -| node[pos=0.25, below, font=\scriptsize, text=thupurple] {运行时反馈 / 新颖性回灌(非梯度迭代,持续开放式进化)} (s1.south);
   % 方法注脚
   \node[font=\scriptsize, text=black!55, below=0.05cm of s1, text width=2.8cm, align=center] {ICM\,/\,RND\\覆盖率缺口};
   \node[font=\scriptsize, text=black!55, below=0.05cm of s2, text width=2.8cm, align=center] {FunSearch\\AlphaEvolve\,/\,EoH};
   \node[font=\scriptsize, text=black!55, below=0.05cm of s3, text width=2.8cm, align=center] {内核 verifier\\隔离沙箱};
   \node[font=\scriptsize, text=black!55, below=0.05cm of s4, text width=2.8cm, align=center] {tail-call\\遥测/回滚};
   \node[font=\scriptsize, text=black!55, below=0.05cm of s5, text width=2.8cm, align=center] {MAP-Elites\\Go-Explore};
  \end{tikzpicture}}
  \caption{防御演化(Defensive Evolution)闭环:以非梯度方式持续进化内核 eBPF 防御。}
  \label{fig:loop}
\end{figure}

\subsection*{5.3 关键方法设计}
\addcontentsline{toc}{subsection}{5.3 关键方法设计}
\begin{itemize}[leftmargin=2em,itemsep=1pt,topsep=2pt]
  \item \textbf{可信闭环(对应关键问题 1)}:四层防御——静态沙箱(禁危险 helper、禁未授权 map 声明、限制循环/规模)$\to$ 模板封装(智能体代码注入安全 wrapper)$\to$ 编译器 `-Wall -Werror -target bpf` $\to$ 内核 verifier;部署后以"丢包异常飙升/健康检查失败"触发自动回滚,关键管理端口(如 SSH)加白名单保护;全部变更落 git 审计。
  \item \textbf{鲁棒探索(对应关键问题 2)}:优先采用对随机性鲁棒的 RND\supercite{burda2019rnd}/伪计数与信息增益信号,而非朴素预测误差;以"防御覆盖率缺口"为显式探索目标。
  \item \textbf{抗腐化进化(对应关键问题 3)}:每条已防住的攻击固化为\textbf{回归用例},智能体每次改动须通过全量攻击回放方可合入;周期性触发"历史压缩"重构,将多条具体特征归并为通用检测原语,以耦合复杂度(coupling complexity)作为系统健康度指标。
  \item \textbf{评测协议(对应关键问题 4)}:构建包含已知攻击与\textbf{预留的未知攻击}的攻防靶场;以攻防协同进化(借鉴 POET 思想)生成递增难度的攻击场景;借助 Cybench/NYU-CTF\supercite{zhang2025cybench,shao2024nyuctf} 思路量化智能体的攻击面理解能力。
\end{itemize}
\subsection*{5.4 实验设计与系统演示}
\addcontentsline{toc}{subsection}{5.4 实验设计与系统演示}
在受控攻防靶场中部署 victim/attacker/puppet/网关与本系统,复现课题组 PMTUD\supercite{feng2025pmtudisolation}、NAT 劫持\supercite{yang2024natwifi}、IPID\supercite{feng2020mixedipid} 等攻击作为红队,度量:① 对已知攻击的检出率/误报率;② 对未知攻击的自主防御提升(迭代曲线);③ XDP 吞吐/时延开销;④ 自动化程度(人工介入量)。与 Kgent\supercite{yang2024kgent}(LLM 写 eBPF)、Tetragon\supercite{tetragon2022}/Falco\supercite{falco2018}(人工策略运行时防御)正面对比。按计算机系答辩要求,提供可测试的典型系统演示(智能体自主从"被某攻击绕过"到"生成有效 eBPF 防御"的端到端闭环)。

\section*{六、可行性分析}
\addcontentsline{toc}{section}{六、可行性分析}
\begin{enumerate}[leftmargin=2em,itemsep=1pt,topsep=2pt]
  \item \textbf{理论可行}:防御演化的方法论根基已被顶刊/顶会验证——FunSearch(Nature)\supercite{romeraparedes2024funsearch}、AlphaEvolve\supercite{novikov2025alphaevolve} 证明 LLM 进化式程序搜索可在系统级代码上做出可验证改进;Darwin–Gödel Machine\supercite{zhang2025dgm} 证明自进化智能体 + 档案库 + 沙箱可行;Kgent\supercite{yang2024kgent} 证明 LLM 能合成正确 eBPF。
  \item \textbf{工程可行}:课题组已有开源系统 \textbf{PacketScope/Guarder} 作为基线,具备 XDP 数据面、cilium/ebpf(bpf2go CO-RE)一键编译、规则/程序热加载、MCP 接入、多维遥测,并已在作者的前期工作中实现"Agent 可编程 eBPF 过滤器 + tail-call 热加载 + 四层安全模型"原型,且通过了 312 个测试用例,\textbf{可行性已被前期原型初步证明}。
  \item \textbf{数据/攻击样本可行}:作者及课题组在 off-path 协议攻击上的系列工作\supercite{feng2025cacm,feng2020mixedipid,feng2022pmtud,feng2025pmtudisolation,yang2024natwifi,wang2025packetsize} 提供了高质量、可复现的真实攻击样本与攻防靶场,是本课题"红队/种子攻击"的独特优势。
  \item \textbf{安全可控}:内核 eBPF verifier 提供了"不可伪造"的形式化安全闸门,配合隔离沙箱与自动回滚,使"让 AI 写内核代码"的风险可控。
  \item \textbf{条件可行}:导师徐恪教授课题组在网络协议安全方向具备深厚积累与算力/实验条件;作者已完成相关攻击研究与系统原型,具备 eBPF、内核与大模型智能体开发能力。
\end{enumerate}
主要风险与对策:① 大模型生成内核代码的稳定性——以 verifier+沙箱+回归测试兜底,必要时引入"浅层模型做感知 + 启发式做规则"的混合;② 自动闭环在真实噪声下未必如靶场般干净——以鲁棒新颖性信号 + 人工门控分级放行缓解;③ 未知攻击的"未知"难以穷尽——以预留测试集 + 协同进化近似,并如实报告覆盖边界。

\section*{七、工作特色与创新点}
\addcontentsline{toc}{section}{七、工作特色与创新点}
\begin{enumerate}[leftmargin=2em,itemsep=1pt,topsep=2pt]
  \item \textbf{可运行的系统(而非设想)}:在 PacketScope/Guarder 上落地一套由编码智能体自动维护的内核 eBPF 防御系统,可演示、可复现。
  \item \textbf{智能体写内核代码的可靠流水线}:把"生成 $\to$ clang 编译 $\to$ 内核 verifier $\to$ 隔离沙箱 $\to$ tail-call 热加载 $\to$ 自动回滚"组织成一条可靠流水线,系统化解决"让 AI 写内核代码"的可用性与安全性(其中 verifier、沙箱、回滚等均为现成基础设施,本文的工作是把它们与编码智能体闭环组织起来,而非重新发明)。
  \item \textbf{面向真实协议攻击的防御评测靶场}:以 PMTUD/NAT/IPID 等真实攻击构建攻防靶场,给出检出率、误报率、性能开销与人工介入量等量化结果,并与人工规则、Kgent\supercite{yang2024kgent}、Tetragon\supercite{tetragon2022}/Falco\supercite{falco2018} 对比;靶场本身可复用于同类内核安全研究。
\end{enumerate}
\section*{八、预期成果}
\addcontentsline{toc}{section}{八、预期成果}
\begin{enumerate}[leftmargin=2em,itemsep=1pt,topsep=2pt]
  \item \textbf{学位论文}一篇,系统阐述防御演化范式、自进化 eBPF 防御系统的设计、实现与评测;
  \item \textbf{开源系统}:在 PacketScope 基础上发布自进化 eBPF 防御系统原型、面向防御演化的\textbf{专用编码智能体 harness}(领域工具 + 闭环编排)与攻防评测靶场;
  \item \textbf{可能的标准化/披露贡献}:若发现并自动缓解新协议攻击,向社区(IETF/Linux 等)负责任披露。
\end{enumerate}
\section*{九、论文工作计划与进度安排}
\addcontentsline{toc}{section}{九、论文工作计划与进度安排}
\begin{table}[htbp]\centering\small
\begin{tabularx}{\textwidth}{l l X X}
\toprule
\textbf{阶段} & \textbf{时间} & \textbf{主要工作} & \textbf{阶段成果} \\
\midrule
开题 & 2026.06 & 选题报告与开题答辩;确定范式与系统框架 & 选题报告、开题答辩 \\
一 & 2026.07–2026.09 & 防御演化范式形式化;在 PacketScope/Guarder 上搭建系统总体框架与闭环骨架 & 系统框架、原型 v0 \\
二 & 2026.10–2026.12 & 新颖性与覆盖度引导的主动探索 + LLM 进化式 eBPF 合成 + verifier/沙箱可信加载闭环 & 闭环原型 v1、初步实验 \\
三 & 2027.01–2027.03 & 多样化防御库维护;攻防协同进化评测;与 Kgent/Tetragon/Falco 对比;\textbf{中期检查} & 完整系统、评测结果、中期报告 \\
四 & 2027.04–2027.05 & 论文撰写、投稿;系统完善;预答辩与盲审送审 & 学位论文初稿、投稿论文 \\
五 & 2027.06 & 论文修改、答辩与学位申请 & 学位论文、答辩 \\
\bottomrule\end{tabularx}\end{table}